\section{A layout for two majority votes}\label{sec:geometry}

Lemma~\ref{lem:vote} delivers a majority symbol to a destination at a
rate bounded by the longest wire from a voter to that destination. We
want two groups of voters to deliver the same symbol to sites \(U,V\).
The resulting shared label has rate \(w(\{U,V\})\). We therefore place
the groups so that the distance from \(U\) to \(V\) exceeds the sum of
their two longest wire lengths. A different split may be needed each
time the recursion changes the active voters. A binary code gives
a small fixed collection of splits from which the active counts alone
select a balanced one.

The same layout holds the two anchor symbols. Its outer faces allow us
to prescribe the two surpluses used in Section~\ref{sec:masks} and then
to enlarge every input group for the last mask in
Section~\ref{sec:recurrence}. We use the normalized mean width
\(w(E)=\delta(E)/\delta(S)\) of Section~\ref{sec:words}.

\begin{lemma}[Layout for two votes]\label{lem:layout}
Let \(r\ge4\), \(q=2\), \(H=\log(2r)\), and
\(d=\lceil12H/\log q\rceil\). There are a finite Euclidean site set
\(S\), distinct satellite sites \(x_1,\ldots,x_{r-2}\), a site \(O\),
and a parameter
\[
        0<\rho<10^{-3},\qquad
        \rho^{-1}\asymp d\sqrt{\log(2dq)}
\]
with the following properties.
\begin{enumerate}
\item \emph{Balanced splits.} There are at most \(dq\) bipartitions
\(C_1,C_2\) of the satellite indices into nonempty classes, each with
destinations \(U,V\in S\) satisfying
\begin{equation}\label{eq:split-gap}
 w(\{U,V\})-
 \max_{i\in C_1}w(\{x_i,U\})-
 \max_{i\in C_2}w(\{x_i,V\})\ \ge\ 4\rho .
\end{equation}
For any subset of \(n\ge2\) active satellites, its counts in the
classes determine a split with at least \(n/(32q)\) active satellites
on each side. Put \(p_0=1/(128q)\). A symbol held by more than
\((1-p_0)n\) active satellites then receives more than three quarters
of the active votes on both sides.
\item \emph{Anchor groups.} All satellites, all destinations and \(O\)
belong to a central polytope of normalized width at most \(10^{-3}\).
For each \(s\in\{\rho,2\rho\}\) there are two disjoint anchor groups
\(A,B\subset S\) outside this polytope, with
\[
         w(A\cup B)=1,\qquad w(A)=w(B)=(1-s)/2 .
\]
Adjoining any central site to either group increases its width by at
least \(0.06\). The group cardinalities are the same for both choices
of \(s\), and
\[
  \ell:=\log\bigl(2(|A|+|B|+2)\bigr)=O(\log(2dq)),
  \qquad \log|S|=O(H).
\]
\item \emph{Enlarged groups.} For \(s=\rho\), each of the two anchors
and the \(r-2\) singleton satellite groups can be enlarged within \(S\)
to increase its width by exactly \(\alpha=\rho/(10r)\). The enlarged
groups are disjoint, their union has width one, and their surplus is
\(9\rho/10\).
\end{enumerate}
All widths use the same set \(S\), and all constants are absolute.
\end{lemma}

\begin{proof}
\emph{Balancing the active votes.}
Assign a word \(c_i\in[q]^d\) to satellite \(i\), choosing the words at
pairwise Hamming distance at least \(d/8\). To see that this is
possible, take independent uniform words. For each pair the probability
of distance less than \(d/8\) is at most
\[
  2^{d h_2(1/8)} q^{d/8-d}\le q^{-d/4}\le (2r)^{-3},
\]
by the binomial sum bound, where \(h_2\) is binary entropy and
\(h_2(1/8)<5/8\). A union bound over the pairs gives the claim.

A coordinate \(j\) and a letter \(b\) specify the class
\(C_1=\{i:c_i(j)=b\}\) and its complement \(C_2\). Keep the splits
with both classes nonempty. Among any \(n\ge2\) active satellites,
each ordered pair of distinct names disagrees on at least \(d/8\)
coordinates. Some coordinate therefore has at least \(n(n-1)/8\)
ordered disagreements. At this coordinate let \(a\) be the largest
active letter count. Each disagreeing pair has at least one member
outside the class of size \(a\), so \(2n(n-a)\ge n(n-1)/8\).
Thus \(a\ge n/q\) and \(n-a\ge(n-1)/16\ge n/32\).
We may choose the first split in a fixed order with at least
\(n/(32q)=4p_0 n\) active votes on each side. Under the majority
promise, fewer than \(p_0n\) votes in the entire active set dissent.
Even if they all lie on one side, they occupy less than one quarter
of its active votes.

\emph{Separating the destinations.}
Use \(m=dq\) orthogonal unit vectors \(e_{j,b}\) and put
\(v_i=\sum_{j=1}^d e_{j,c_i(j)}\). For the split \((j,b)\), write
\(e=e_{j,b}\). Every vector has squared norm \(d\), but its projection
onto \(e\) is one for \(C_1\) and zero for \(C_2\). Place the
destinations at \(Te,-Te\), respectively, with \(T=10d\).
Figure~\ref{fig:vote-layout} shows the two routes. The projection of
\(C_1\) toward its destination reduces the combined length of the
longest wires. In unscaled coordinates their lengths sum to
\[
   \sqrt{T^2-2T+d}+\sqrt{T^2+d}
      \le 2T-1+d/T=2T-9/10 .
\]
Here we used the tangent-line bound
\(\sqrt{T^2+y}\le T+y/(2T)\). The distance between the destinations
is \(2T\), leaving an unscaled gap of at least \(9/10\).

\begin{figure}[!htbp]
\centering
\begin{tikzpicture}[x=1cm,y=1cm,font=\small,
    wire/.style={draw=sea,thick},shadow/.style={draw=ink!55,densely dashed}]
  \node[font=\small\bfseries] at (0,1.68) {A split along \(e_{j,b}\)};
  \draw[draw=ink!70,->] (-2.44,0)--(2.7,0);
  \draw[shadow] (0,-.92)--(0,.92);
  \draw[shadow] (.85,-.92)--(.85,.92);
  \node at (0,1.16) {\(C_2\)};
  \node at (.85,1.16) {\(C_1\)};
  \foreach \y in {-.69,-.25,.46}{\fill[ink] (0,\y) circle (1.6pt);}
  \foreach \y in {-.48,.34,.72}{\fill[ink] (.85,\y) circle (1.6pt);}
  \draw[wire] (-2.2,0)--(0,-.69) (.85,.72)--(2.2,0);
  \fill[sea] (-2.2,0) circle (2pt) node[below=3pt,text=ink] {\(V=-Te\)};
  \fill[sea] (2.2,0) circle (2pt) node[below=3pt,text=ink] {\(U=Te\)};
  \node at (0,-1.12) {\(0\)};
  \node at (.85,-1.12) {\(1\)};
  \begin{scope}[xshift=6.2cm]
    \node[font=\small\bfseries] at (0,1.68) {The two anchor groups};
    \fill[sea!14] (-1.85,1) rectangle (1.85,1.24);
    \fill[sea!14] (-1.85,-1.24) rectangle (1.85,-1);
    \draw[draw=ink] (-1.85,-1.24) rectangle (1.85,1.24);
    \draw[draw=ink!65] (-1.85,1)--(1.85,1)
                       (-1.85,-1)--(1.85,-1);
    \draw[draw=rust,thick,dashed] (-1.85,1.24)--(0,0)--(1.85,1.24);
    \fill[ink] (0,0) circle (1.8pt) node[below] {\(z\)};
    \node[anchor=west] at (1.94,1.12) {\(A\)};
    \node[anchor=west] at (1.94,-1.12) {\(B\)};
  \end{scope}
\end{tikzpicture}
\caption{At left, the two classes project to zero and one on the chosen
code axis. The highlighted wires carry their votes to opposite
destinations. At right, a projection of the prism onto its interval
and height coordinates shows both anchor layers. Adding a central
site \(z\) to the upper anchor places the dashed triangle in the
projection of its hull. Distances and thicknesses are schematic.}
\label{fig:vote-layout}
\end{figure}

All \(v_i\) and all destinations lie in the cross-polytope
\(T B_1^m\), which has only \(2m\) vertices. For independent standard Gaussian coordinates \(g_j\), its mean width satisfies
\[
 \delta(TB_1^m)=\sqrt{2\pi}\,T\,\E\max_{j\le m}|g_j|
       \le \sqrt{2\pi}\,T\sqrt{2\ln(2m)} .
\]
Indeed, for \(u>0\) the expected
exponential of \(u\max_j|g_j|\) is at most \(2m e^{u^2/2}\).
Taking logarithms and minimizing the resulting upper bound on the
mean gives the inequality. Scale all positions by
\(\gamma=10^{-5}/(d\sqrt{\log(2m)})\). The resulting polytope
\(K=\gamma T B_1^m\) has \(\delta(K)\le10^{-3}\), and the wire gap
is at least \(0.9\gamma\).

\emph{Setting the anchor widths.}
Take the prism
\[
       P=P_0\times[-h,h],\qquad P_0=[-1,1]\times K,\qquad h=1.025
\]
in orthogonal coordinates. The central polytope consists of points
\((0,x,0)\) with \(x\in K\). Place the satellites at
\((0,\gamma v_i,0)\), the destinations at \((0,\pm\gamma Te,0)\),
and \(O\) at the origin. The upper anchor \(A\) contains all vertices
of \(P_0\) at height \(h\), and the lower anchor \(B\) contains their
reflections at \(-h\). Together these vertices span \(P\). We will
add sites only inside the prism, fixing the normalizing denominator
\[
       D=\delta(S)=\delta(P)=\delta(P_0)+2h,\qquad
       2\le\delta(P_0)=2+\delta(K)\le2.001 .
\]
In particular \(4.05\le D\le4.051\). Set \(\rho=\gamma/32\).
After normalization the wire gap is at least \(0.9\gamma/D>4\rho\),
proving~\eqref{eq:split-gap}, and the central polytope has width at
most \(10^{-3}/D\).

For each \(s=\rho,2\rho\), include in each anchor a second copy of
its vertices, a distance \(t_s\) toward the center. Its hull is
a product of \(P_0\) with an interval of length \(t_s\), giving width
\((\delta(P_0)+t_s)/D\). The choice
\[
       t_s=\tfrac12(1-s)D-\delta(P_0)
           =\tfrac12\bigl(2h-\delta(P_0)-sD\bigr)
\]
lies between \(0\) and \(0.03\) and gives the desired normalized
width \((1-s)/2\). Both choices use the same outer faces.

If \(z=(0,x,0)\) is central, the hull of \(A\cup\{z\}\) contains
the triangle joining \(z\) to \((-1,0,h)\) and \((1,0,h)\) in the
outer face. Its raw width is its semiperimeter,
\(1+\sqrt{1+h^2+\|x\|^2}>2.4\). Thus \(w(A\cup\{z\})>2.4/D>0.56\),
whereas \(w(A)\le1/2\). Reflection gives the same gain for \(B\).

For \(s=\rho\), append to each anchor one more layer a distance
\(\alpha D\) farther inward and to each satellite one more site at
height \(\alpha D\) above it. Each anchor hull grows by an interval
of length \(\alpha D\), and each satellite group becomes a segment of
that length. Hence each of the \(r\) widths increases by exactly
\(\alpha\). The anchor layers stay within \(0.04\) of their respective
outer faces, and different satellite groups have different projections
onto \(K\). The groups remain disjoint, and their union still spans
\(P\). Their surplus is therefore
\[
                  1-\bigl(1-\rho+r\alpha\bigr)=9\rho/10 .
\]
Each layer has \(4m\) vertices. Each original anchor has two layers,
and only a constant number of layers is needed altogether. Including
the satellites, their extra sites, the destinations and \(O\) gives
\(|S|=O(r+dq)\) and the stated bound on \(\ell\).
\end{proof}
